\begin{pastexam}[label={px:df3-114-ds-04}]{地方三等資訊處理 114 年 第四題}
（申論；25 分）

一個自動化工廠大量採用機器人協助裝箱作業。該工廠固定時間生產出一組 $n$ 個不同大小的塑膠球並放到裝箱作業
輸送帶上。輸送帶上配置數個機器人，當輸送帶上的球經過時，機器人負責將眼前兩顆球將順序排序正確，大的在前，
小的在後。當輸送帶上的球經過了所有機器人後，球的順序就完全由大排到小了。（每小題 5 分，共 25 分）
\begin{enumerate}[leftmargin=2.2em]
\item[㈠] 若 $n = 6$，且生產後放上裝箱輸送帶的球的大小為 3, 2, 5, 6, 1, 4。請說明若輸送帶配有 4 個機器人
是否足夠將球的順序完全由大排到小？
\item[㈡] 若 $n = 20$，且生產後放上裝箱輸送帶的球的大小為 11, 12, 20, 16, 3, 1, 7, 15, 2, 18, 10, 5, 14, 6, 8,
13, 19, 4, 9, 17，請說明輸送帶上最少該配置幾個機器人才能將球的順序由大排到小？
\item[㈢] 若 $n = 6$，且輸送帶上配有 4 個機器人，請給一組放上裝箱輸送帶的球的大小順序，使得其經過這 4 個
機器人後，整組球的順序仍未能排好。
\item[㈣] 若每一組球生產後放上裝箱輸送帶的\underline{球的大小順序非固定順序}，請說明輸送帶上最少該配置幾個
機器人才能每次都能將球的順序由大排到小？
\item[㈤] 若 $n = 10$，且每一組球生產後放上裝箱輸送帶的\underline{球的大小順序非固定順序}。假設輸送帶上原本
配置 $n$ 個機器人，若改成配置 $2n$ 個機器人，整組球順序排好的速度可以加快多少？請說明。
\end{enumerate}
\end{pastexam}

\begin{pastsoln}
\textbf{模型。} 球依序經過一個機器人時，機器人手上留著一顆球，下一顆到的時候比較兩顆，把大的放
出去（往前），留下小的，再和下一顆比。 整組球經過一個機器人，等於對序列做\textbf{一輪泡沫排序}：
由前往後比較相鄰兩顆，前面比較小就交換，一輪下來最小的球被帶到最後面（資結書第~14.2~節，這裡是
由大到小排）。 經過 $k$ 個機器人就是做 $k$ 輪。 題目沒有更精確地描述機器人的動作，以下都用這個模型。

\textbf{需要幾輪。} 一輪之中，一顆球最多往前移一格（它只會在和前一顆交換時往前），而只要前面還有比它
小的球，這一輪它就一定往前移一格（帶著前面那些較小球之一往後走的那顆球會經過它）。 由大到小排時，
一顆球最後要排在所有比它小的球前面。 所以令
\[
c(x) = \text{排在 } x \text{ 前面、而且比 } x \text{ 小的球的個數},
\]
需要的輪數（機器人數）就是 $\max_x c(x)$。

\textbf{㈠ 夠，只需要 3 個。} 3, 2, 5, 6, 1, 4 的 $c$ 值依序是 0, 0, 2, 3, 0, 3，最大是 3。 實際逐輪：
\begin{center}
\begin{tabular}{cl}
\toprule
經過 & 順序 \\
\midrule
開始 & 3, 2, 5, 6, 1, 4 \\
第 1 個 & 3, 5, 6, 2, 4, 1 \\
第 2 個 & 5, 6, 3, 4, 2, 1 \\
第 3 個 & 6, 5, 4, 3, 2, 1 \\
第 4 個 & 6, 5, 4, 3, 2, 1（沒有交換） \\
\bottomrule
\end{tabular}
\end{center}
3 個機器人之後已經由大排到小，4 個足夠。

\textbf{㈡ 最少 16 個。} 對每顆球數前面比它小的個數：
\begin{center}
\small
\begin{tabular}{l*{20}{c}}
\toprule
球 & 11 & 12 & 20 & 16 & 3 & 1 & 7 & 15 & 2 & 18 & 10 & 5 & 14 & 6 & 8 & 13 & 19 & 4 & 9 & 17 \\
$c$ & 0 & 1 & 2 & 2 & 0 & 0 & 2 & 5 & 1 & 8 & 4 & 3 & 8 & 4 & 6 & 10 & 15 & 3 & 8 & 16 \\
\bottomrule
\end{tabular}
\end{center}
最後一顆 17 前面有 16 顆比它小（除了 20、18、19 以外都是），它要往前移 16 格才到正確位置（第 4 名，
前面只有 20、19、18），每經過一個機器人最多前進一格，所以至少要 16 個；由上面的論證 16 個也夠。
答案 16 個。

\textbf{㈢ 例如 1, 2, 3, 4, 5, 6。} 6 前面有 5 顆比它小的球，每經過一個機器人只前進一格，經過 4 個後
還在第 2 個位置：逐輪是 2,3,4,5,6,1 → 3,4,5,6,2,1 → 4,5,6,3,2,1 → 5,6,4,3,2,1，最後是
5, 6, 4, 3, 2, 1，還沒排好。 任何讓某顆球的 $c$ 值是 5 的順序都可以，例如 6 排在最後、前面五顆任意排（如 5, 4, 3, 2, 1, 6）。

\textbf{㈣ $n - 1$ 個。} 一顆球的 $c$ 值最多 $n - 1$（前面全部都比它小），所以 $n - 1$ 個機器人一定夠。
也一定要這麼多：順序不固定時，可能遇到最大的球排在最後面（例如由小到大的 $1, 2, \dots, n$），它要往前
移 $n - 1$ 格，少一個機器人就排不好。 （等價的說法：每經過一個機器人，至少有一顆球確定排到最後的
正確位置，$n - 1$ 個之後剩下的一顆自然也在正確位置。）

\textbf{㈤ 不會加快。} 由㈣，$n = 10$ 時 $9$ 個機器人就保證排好，原本配置 $n = 10$ 個已經足夠，排好所需的
機器人數（輪數）由球的順序決定，最多 9。 改成 $2n = 20$ 個，多出來的 10 個機器人經過時球已經排好，
只做比較、不會交換，對排序沒有幫助，排好的時間點一樣，加快 0 倍；如果每個機器人處理一組球要花固定
時間，球還要多經過 10 個機器人，整組出輸送帶反而更晚。 要真正加快，得改變機器人的工作方式（例如
讓多個機器人同時處理不同位置的相鄰兩顆球，像奇偶轉置排序那樣平行比較），單純增加同樣工作的
機器人數沒有用。
\end{pastsoln}
